\documentclass[10pt,twocolumn]{article}
\usepackage[margin=0.72in]{geometry}
\usepackage{amsmath,amssymb,amsthm,booktabs,array,longtable}
\usepackage[T1]{fontenc}
\usepackage{lmodern}
\usepackage[hidelinks,breaklinks]{hyperref}
\usepackage{microtype}
\usepackage{balance}
\setlength{\columnsep}{0.25in}
\setlength{\emergencystretch}{1em}
\newcolumntype{P}[1]{>{\raggedright\arraybackslash}p{#1}}
\newtheorem{theorem}{Theorem}
\newtheorem{proposition}[theorem]{Proposition}
\newtheorem{corollary}[theorem]{Corollary}
\newtheorem{definition}{Definition}
\newcommand{\src}[1]{\nolinkurl{#1}}
\title{Analyzing Delegation and Handoff Locks:\\Separate Performance and Fairness Models}
\author{Research section draft}
\date{}
\begin{document}
\maketitle
\begin{abstract}
We give separate analyses of lock performance and service fairness.
The performance analysis uses a LogP-inspired communication and execution
model to derive migration amortization, throughput crossover, and
caller-return latency conditions for a given service policy and workload.
The fairness analysis independently defines service credit and eligibility,
then derives service-spread, accounting-error, and competing-service bounds
without communication-cost assumptions. Separate implementation comparisons
show why queue-pop budgets differ from completion budgets, equal billed
service need not mean equal useful service, and minimum-credit ordering
alone does not establish a dynamic-arrival guarantee. A final discussion
states the additional assumptions needed to relate the two analyses.
The mathematical artifact checks finite cases and explicit counterexamples;
it establishes neither calibrated hardware performance nor concurrency
correctness of the implementations.
\end{abstract}
\section{Scope and analytical contract}
\label{sec:scope}
We analyze performance and fairness as two separate questions.
Flat combining~\cite{fc} and queue-based combining~\cite{cc} decouple the
requester from the executor; handoff locks normally execute the critical
section on its owner. The same implementation can have a favorable execution
cost and an inadequate service guarantee, or the reverse.

Section~\ref{sec:performance} models communication, execution placement,
scheduler work, throughput, and return latency for a given service policy.
Section~\ref{sec:fairness} defines the service objective and proves bounds
from selection, visibility, admission, and accounting assumptions, without
LogP parameters. Each analysis has its own implementation comparison and
validation criteria. Only Section~\ref{sec:discussion} relates their outputs.
Batch amortization and minimum-credit scheduling are not claimed as novel.
We supply conditional results and source evidence, not memory-safety proofs,
calibrated hardware predictions, or empirical validation from finite checks.

The primary scope is all 21 DLock2 targets at commit \texttt{56a02ab};
Appendix~\ref{app:inventory} also maps all 11 DLock1 targets and dispatch
limitations. We assume non-recursive operations on one protected object,
one outstanding synchronous call per requester, and finite non-preemptive
delegates. Performance estimates initially assume pinned, non-oversubscribed
workers; fairness uses explicitly stated cohort and selection assumptions.
Blocking, preemption, external accesses, and stronger progress or residency
conditions are not silently treated as harmless.

\begin{definition}[Request and execution epoch]
A request has events $a$ (invocation), $p$ (publication), $v$ (visibility to the
service selector), $s$ (service start), $f$ (service finish), $n$ (completion
notification), and $r$ (caller return). These are causally ordered, but need
not be adjacent. An execution epoch is one successful combining pass/batch
or one handoff critical section. Consecutive epochs may use the same core.
\end{definition}
Visibility is policy-specific: a node in an announcement ring or a temporary
completed-node buffer is not necessarily in the current selectable set.
Completion is not return: a caller acting as combiner may execute more
requests after its own result is available.

\paragraph{A concrete correctness prerequisite.}
FC-SL's callers read the global non-atomic usage/served counters before
acquiring the combiner lock, while a combiner writes them. This is a
source-level data-race blocker, not a bounded timing-error term
(Appendix~\ref{app:fcsl-race}). Its abstract scheduler can be analyzed, but
implementation-level fairness/performance guarantees require synchronization
repair first. We retain the source unchanged rather than hide the issue by
analyzing an implicitly repaired implementation.

% Performance model and throughput results.
\section{Performance analysis}
\label{sec:performance}
This analysis asks how much communication and execution work each lock
requires, and how that work constrains throughput and caller-return latency.
The service policy, executed operation mix, and executor sequence are inputs;
no equal-share or bounded-lag guarantee is assumed. Results here do not use
the fairness theorems. Costs are elapsed time, resource occupancy, or source
work counts, not a scheduler's service-credit metric.

\subsection{Communication and execution model}
\label{sec:model}
\subsubsection{Communication events}
Classical LogP~\cite{logp} specifies small-message latency bound $L$, processor
send/receive occupancy $o$, endpoint send/receive gap $g$, and $P$ processor
modules. Its capacity is at most $\lceil L/g\rceil$ messages in flight from
or to an endpoint. The runtime convention charges $L$ per message, giving
$L+2o$ for an isolated one-way delivery. Independent deliveries can overlap:
latency belongs to dependency paths, while occupancy and gaps constrain rate.
An upper-bound $L$ is not a physical lower-bound latency.

For coherent shared memory we use a \emph{LogP-inspired adaptation}, not an
identification of a load with a message. Repeated private-cache polling can
produce no new fabric traffic. Ownership transfer can require requests,
invalidations, acknowledgments, and forwarding. Following the motivation of
cache-coherent communication models~\cite{coherence}, transaction costs are
indexed by type and topology. $P$ counts hardware execution endpoints and
$N$ counts clients; directory/interconnect bottlenecks may be additional
resources. A large cache line is not automatically a LogGP long message:
LogGP's byte-gap parameter $G$~\cite{loggp} is appropriate only when a calibrated
bulk-transfer regime exists.


Construct a finite event DAG containing delegate execution, publication,
discovery, queue manipulation, transfer, completion, and executor changes.
Delegate nodes touching the same protected state are serialized. Independent
publication and notification may overlap this chain. Label CPU work with its
endpoint, messages with endpoint/topology, and contended atomic ownership
with the resource it occupies. Failed elections, empty scans, and parked
threads have events but contribute no successful operations.

\begin{proposition}[Resource and dependency bound]
\label{prop:resource}
For a closed finite DAG completing $B$ operations, let $\mathrm{CP}$ be its
longest weighted path, $W_p$ its total non-overlappable CPU occupancy at
endpoint $p$, $v_p^s,v_p^r$ its message-send/receive counts, and $r_\ell a_\ell$
the occupancy of a serialized hot-line resource. With endpoint gap $g_p$,
\begin{align}
T &\ge \max\{\mathrm{CP},\max_p W_p,\max_p G_p,
                    \max_\ell r_\ell a_\ell\},\label{eq:resource}\\
G_p&=\max\{(v_p^s-1)_+,(v_p^r-1)_+\}g_p .\nonumber
\end{align}
The costs must be exact within the abstract machine or guaranteed minima.
Capacity and any additional fabric constraints remain part of the DAG/model.
\end{proposition}
\begin{proof}
Every dependency path must complete; each non-overlappable resource must
supply its total occupancy. A sequence of $v$ sends or receives has $v-1$
inter-event gaps. Each requirement separately lower-bounds elapsed time;
their maximum is therefore a lower bound. Taking a sum would unjustifiably
forbid overlap. Include startup/drain events rather than counting transfers
that lie outside the chosen observation interval.
\end{proof}
Thus $B/T$ has a model upper bound, not necessarily an attainable rate.
Substituting measured \emph{means} for guaranteed minima yields a performance
estimate, not a hardware theorem. Moreover, $O(1)$ successful queue operations
do not imply $O(1)$ elapsed time when a predecessor can stop publishing.

\subsubsection{Trace accounting and a predictive approximation}
Let a trace have $E$ passes, $B>0$ completions, and $F$ executor changes.
Define $b=B/E$ and $p=F/E$; zero-completion passes count in $E$.
Partition the \emph{serialized execution lane}, without double counting,
into warm service, per-operation scheduling/discovery and exposed request
communication, pass administration, exposed protected-state/scheduler-state
reacquisition, and uncovered idle time. Mean per-completion time is then
represented by
\begin{equation}
 t_D \simeq \bar C_D+d_D+
 \frac{A_D+p_D(M_D+K_D)}{b_D}+I_D .\label{eq:delegation}
\end{equation}
$A$ is mean administration per pass, $M$ and $K$ are mean exposed protected-data
and scheduler-state reacquisition \emph{conditional on an executor change},
and $I$ is uncovered idle time per completion. When $F=0$, its migration term
is zero; finite startup cost is separate. $\bar C$ and $d$ are weighted by
operations, not passes. Ratios of totals give this accounting; averaging
per-pass ratios or assuming independence of batch size and cost does not.

As a trace partition, the accounting is definitional. As a prediction,
Equation~\eqref{eq:delegation} replaces trace costs by independently calibrated
values and is only an approximation. It does not replace the DAG constraints.
A request miss inside measured service must not be added again to $d$;
a cold protected-data miss must not be counted both in $C$ and $M$.
Speculative/prefetched communication contributes only its \emph{exposed}
residual to this serialized decomposition. Network demand still appears in
Equation~\eqref{eq:resource} even if its latency was hidden.

Protected-data cost $M(W,\tau,\rho)$ depends on touched reusable lines $W$,
topology $\tau$, and residency $\rho$, not merely allocated object size.
Independent transfers might fit $d_\tau+(W-1)\gamma_\tau$; a dependent
pointer chain might fit $W d_\tau$. These are alternative calibrated regimes,
not universal coherence formulas. Input lines, queue metadata and protected
state are separate working sets. A large heap can cease to be cache resident.
An unchanged executor can also suffer capacity misses or interference;
those costs are not excluded by the absence of executor-change migration.

A comparable handoff estimate is
\begin{equation}
 t_H\simeq \bar C_H+h_H+M_H,\label{eq:handoff}
\end{equation}
where $h_H$ includes lock protocol and policy work, and $M_H$ is actual exposed
movement per operation, possibly zero for local reuse. Admission cooldowns
are not automatically $I_D$: another request may use the server meanwhile.

\subsection{Migration amortization and throughput crossover}
\label{sec:amortization}
\begin{proposition}[Intra-epoch ownership invariance]
\label{prop:ownership}
Fix an epoch's executor, its operations' accessed-line trace up to a
permutation, and exclusive protected-state access by that executor during the
epoch. Permuting request order cannot introduce a transfer of a protected
line to a \emph{different executing core within that epoch}.
\end{proposition}
\begin{proof}
Every protected access in the permuted epoch is still issued by the same
core. There is no second executor to which execution ownership can transfer.
This says nothing about first-touch acquisition from an earlier owner,
eviction/refetch from memory, requester inputs, or off-path metadata traffic.
\end{proof}
This is narrower than ``the data always stays in L1.'' If ordering changes
which lines are accessed, the initial owners of those lines, the operation
mix, batch boundaries, or executor selection, even the migration comparison
needs to be recomputed.

\begin{corollary}[Resident repeated-footprint case]
For a fixed $W$-line exclusively modified footprint touched in every epoch,
no external invalidations/evictions, and $F$ executor changes, at most $WF$
protected-line ownership changes are required after warmup; under the model
where every changed executor must acquire every line, the count is exactly
$WF$. Per completion it is $Wp/b$.
\end{corollary}
\begin{proof}
An unchanged executor requires no reacquisition. Each changed executor
acquires each of the $W$ lines once, then retains it for that epoch.
Sum over changed epochs and divide by $B$, using $F/B=p/b$.
\end{proof}
Handoff locks correspond to one operation per epoch, but their switch
probability and topology must be measured. Owner-to-owner movement already
exists in MCS; it is not a cost introduced only by CFL's added policy.

\begin{proposition}[Conditional batching crossover]
\label{prop:crossover}
Assume the same service mix and warm cost $C$, the same migration cost $M$,
$p_D=1$, no uncovered idle time, and costs independent of batch size in the
considered regime. Within Equations~\eqref{eq:delegation}--\eqref{eq:handoff},
delegation has smaller time per operation exactly when
\begin{equation}
 b(M+h-d)>M+A+K .\label{eq:crossover}
\end{equation}
If $M+h-d>0$, the smallest positive integer completion batch satisfying this
inequality is
\begin{equation}
 b_{\min}=\left\lfloor\frac{M+A+K}{M+h-d}\right\rfloor+1.
\end{equation}
If the denominator is nonpositive and all costs are nonnegative, no positive
batch meets the strict inequality.
\end{proposition}
\begin{proof}
Cancel $C$ in $C+d+(M+A+K)/b<C+h+M$ and multiply by $b>0$.
The integer result follows from the strict inequality.
\end{proof}
The saved term is $(1-1/b)M$; the unpaid terms are per-request scheduling and
communication plus $(A+K)/b$. Increased migration cost lowers the relative
importance of scheduler work, but oversized inputs, a congested publication
line, or migrating PQ metadata can remove the advantage. This is a
regime-dependent prediction, not a ranking of all locks on all machines.

\subsubsection{Implementation deltas and controlled operation mix}
To explain the performance difference between two variants, compare:
\begin{align}
\Delta_D&\simeq\Delta \bar C+\Delta d+
 \Delta\!\left(\frac{A+p(M+K)}b\right)+\Delta I,\\
\Delta_H&\simeq\Delta\bar C+\Delta h+\Delta M.
\end{align}
FC-PQ versus FC changes discovery and batching as well as policy; CC-Ban
versus CC also changes the source budget (16 versus 64). CFL versus MCS
changes queue policy and locality, not the existence of handoff itself.
None is automatically a single-variable ablation. The evidence appendix
identifies these structural confounders before interpreting a slowdown.

For fixed per-class service $c_i>0$ and given useful-service fractions $f_i$
on a fully utilized overhead-free server,
\begin{equation}
 X_i=f_i/c_i,\qquad X=\sum_i f_i/c_i.\label{eq:mix}
\end{equation}
For costs 1 and 9, the supplied mix $(f_1,f_2)=(1/10,9/10)$ gives
$X=1/5$, whereas $(1/2,1/2)$ gives $X=5/9$ operations per normalized time
unit. No overhead changed. These fractions are inputs to the cost model,
not conclusions about which shares an implemented scheduler guarantees.
Report useful work, class completion rates, service shares, and total
operations together. A change in $\bar C$ due to a different class mix is not
necessarily synchronization inefficiency.

% Performance latency and implementation costs.
\subsection{Batching versus caller return latency}
\label{sec:return}
A waiter can return after its own notification, while an executor caller may
finish the rest of the batch first. If its operation is first in a batch of
exactly $b$ completed operations, each later operation with serialized elapsed
service at least $\tau_{\min}$ adds at least $(b-1)\tau_{\min}$
post-completion delay. This is an elapsed-time lower bound, excluding extra
administration. It follows from the synchronous return rule regardless of
the policy that selected the operation order.

Suppose a restricted implementation/regime has a certified per-subsequent-
operation serialized upper bound $\tau_{\max}$ and remaining fixed tail work
$e_{\rm tail}$. Then
\begin{equation}
 R_{\rm tail}\le e_{\rm tail}+(b-1)\tau_{\max}.\label{eq:tail}
\end{equation}
This assumes exactly $b$ completed operations and that \emph{all} other
post-completion work is in the envelope. It is not a bound for FC-PQ obtained
by replacing $b$ with 64: its announcement work, no-service pops, stalled
publishers, and retries also need bounds.

\begin{proposition}[Certified batching feasibility]
\label{prop:frontier}
Assume $\tau_{\max}>0$. Under
Proposition~\ref{prop:crossover}'s performance approximation and the
return envelope~\eqref{eq:tail}, a tail budget $D\ge e_{\rm tail}$ admits a
positive integer batch certified by both models iff
\begin{equation}
 b_{\min}\le b_{\max}:=
 1+\left\lfloor\frac{D-e_{\rm tail}}{\tau_{\max}}\right\rfloor,
 \qquad \tau_{\max}>0,
\end{equation}
with positive crossover denominator. If either condition fails, no batch
has both certificates under the assumed envelopes.
\end{proposition}
\begin{proof}
The strict throughput inequality requires $b\ge b_{\min}$.
Rearranging~\eqref{eq:tail} to meet $D$ gives $b\le b_{\max}$.
The integer interval is nonempty exactly under the stated condition.
\end{proof}
For an implementation with supported batches $\mathcal B$, further intersect
this integer interval with $\mathcal B$ and with the regime in which the
cost/envelope assumptions hold. A zero $\tau_{\max}$ is outside the stated
proposition; it would impose no batch-dependent tail limit when
$D\ge e_{\rm tail}$.
Absence of a certificate is not proof that real hardware cannot meet both
objectives: an upper envelope may be conservative. As a synthetic example,
let $(M,A,K,h,d)=(20,4,4,2,8)$, giving $b_{\min}=3$.
For $(e_{\rm tail},\tau_{\max},D)=(3,5,13)$, $b_{\max}=3$;
only $b=3$ is certified. Reducing $D$ to 8 gives $b_{\max}=2$ and no
certified batch. All numbers are normalized model costs, not measurements.
This locates the point where merely increasing a batch budget cannot satisfy
both modeled aims; a different return/executor mechanism or smaller
coordination cost is needed to improve the certified region.

For a stationary closed workload, the operation-weighted response-time law
$N=X(\mathbb E[R]+\mathbb E[Z])$, with $Z$ the outside-lock think time,
is an aggregate consistency check. It does not predict percentiles or an
individual request's maximum waiting time.

\subsection{Performance of the implemented lock families}
\label{sec:instances}
Table~\ref{tab:performance} summarizes the dominant \emph{source work}, not exact
fabric traffic. The appendix provides complete target coverage and source
ranges. Let $n$ be retained queue/list size, $a$ arrivals drained, $k$ pop
attempts, and $b$ actual completions; use $n_{\max}$ over a pass for complexity
bounds. Additional failed elections, allocation, stalls, and cleanup are
not assumed constant.

\begin{table*}[t]
\centering\small
\caption{Performance comparison only. Complexity describes source work,
not latency or one-message-per-instruction counts. Elapsed-time predictions
add machine, residency, and progress assumptions; no fairness theorem is used.}
\label{tab:performance}
\begin{tabular}{@{}P{0.14\textwidth}P{0.28\textwidth}P{0.25\textwidth}P{0.24\textwidth}@{}}
\toprule
Family & Discovery/scheduler demand & Execution and budget & Performance consequence \\
\midrule
FC & $O(n)$ visits/pass; $O(n/b)$ visits/completion for $b>0$ & Full list scan, no count cap; executor may change each pass & Retained idle nodes increase scan work; occupancy determines amortization. \\
FC-Ban & FC scan plus eligibility tests, skips and retained-registration accounting & Skips reduce $b$; uncovered idle differs from client cooldown & Charge scan visits and exposed idle, not every caller's entire cooldown. \\
FC-PQ heap/tree & $O((a+k)\log(n_{\max}+1))$ ordinary/amortized queue work plus buffers/timers & $a\le64$, $b\le k\le64$ per pass; queue may exceed 64 clients & Activation/pop ratios and migrating PQ metadata determine scheduling cost. \\
FC-SL & Ordered skip-set operations, allocation and atomic metadata costs & Count budget on pops, not completions; direct ordered publication & Shared ordered metadata adds traffic; current source has the shared correctness blocker. \\
CC / DSM & Linked-queue publication and $O(b)$ service traversal & CC cap 64, DSM source permits 65; both may stop earlier & Batch occupancy amortizes executor changes; absent links can stall progress. \\
CC-Ban & CC protocol plus pre-enqueue wait and charge accounting & Cap 16, unlike CC's 64 & Comparative cost includes budget difference and admission effects. \\
MCS / CLH & Constant successful queue/link operations; local/per-predecessor polling & Caller executes one operation per ownership epoch & Local polling does not remove protected-data migration or delayed linking. \\
Spin / Ticket & Central acquisition/turn cache lines; retry or sharer traffic & Caller execution; no delegation amortization & Contended line demand and placement matter beyond successful atomic counts. \\
Shfl / CFL & Variable queue inspection/relinking; CFL adds usage counters & Locality/policy can alter executor topology and run length & Model policy work and changed migration separately; neither is a fixed universal penalty. \\
Mutex / U-SCL & OS waiting or slice admission plus ownership protocol & Caller execution; kernel/preemption terms may dominate & Runtime implementation and scheduling assumptions are required for latency bounds. \\
\bottomrule
\end{tabular}
\end{table*}

\paragraph{FC versus FC-PQ.}
Under a full active FC scan, $b\approx n$ makes per-operation scan work
approximately constant. Under sparse activity, $n/b$ grows. In FC-PQ,
$k/b$ quantifies wasted completed-node pops, and $a/b$ activation intensity;
neither is fixed by $H$. Its scheduler ratio is approximately
$(a+k)\log(n_{\max}+1)/b$, not unconditionally $O(N\log N)$ per pass.
A heap's growth can copy $O(n)$ storage on a particular insertion, so an
amortized cost is not a per-pass worst-case latency guarantee. BinaryHeap
and BTreeSet share the ordering policy, not identical locality/allocation
costs; the shared PQ migrates when the combiner changes.

\paragraph{Queue and handoff families.}
CC/DSM amortize linked-queue service over an executor epoch. MCS/CLH
reduce centralized polling traffic relative to globally shared-turn polling,
but still have enqueue hotspots and owner-to-owner protected-state movement.
ShflLock and CFL change the topology and order of owners; the model therefore
changes both $M_H$ and policy work. Rust/C wrappers must be treated as
different implementations:
the C FC cleanup/polling rules and C CC batch limit are recorded in the appendix.


\paragraph{Admission and waiting strategies.}
FC-Ban scans and skips published requests; CC-Ban delays the requester before
enqueue. These paths change scan work, actual occupancy, and exposed idle
in different ways. U-SCL's slice admission needs its C protocol, not the
Rust wrapper alone. Mutex and blocking-Parker paths require wakeup and
scheduler terms; they are not accurately ranked by adding a fixed LogP
message latency to a spin-lock cost. A dedicated RCL executor, when enabled,
changes executor migration but consumes a server resource and scan capacity.
The checked-in dispatch limitations matter before claiming such a comparison.

\subsection{Performance validation and limits}
\label{sec:performance-validation}
A quantitative use of the model requires independently calibrated
transaction/topology costs and observable event counts. At minimum record
per-pass $(a,k,b,n)$, scan visits, activation frequency, executor transitions,
class-specific completions and service costs, and full pass/return intervals.
Do not fit a fresh unexplained residual for every lock. The present FC-PQ
\texttt{combiner\_time\_stat} reuses its start timestamp inside each delegate,
so it is not a full-pass timer; the benchmark's requester-equals-executor
flag is not a complete history of who performed combining work.

Two falsifiable regime predictions follow: at matched work and topology,
exposed per-operation migration decreases with $F/B=p/b$ until another
bottleneck dominates; sparse FC activity tracks retained-list visits per
completion, while FC-PQ tracks activation/pop work per completion rather than
only $N$. Neither prediction assumes a service-fairness bound.

Hold out footprints, operation mixes, and placements after calibration, then
predict time and crossover direction. Report uncertainty and failures,
including locality effects outside the repeated-footprint assumption.
Platform-specific ownership events and controlled access patterns are needed
to attribute coherence; generic LLC misses cannot identify protected-state
versus publication/PQ traffic. No coefficients or machine measurements are
fabricated here. The \texttt{crossover} and \texttt{batch\_frontier} artifact
checks exercise the algebra, not hardware speed or individual starvation.

\clearpage
% Fairness metrics, proofs, and implementation obligations.
\section{Fairness analysis}
\label{sec:fairness}
This analysis asks which service metric is equalized, among which clients,
and with what bound on service disparity or competing selections. Its inputs
are the selectable set, admission/initialization rules, the selection rule,
and service charges. It uses no LogP parameters, cache-residency assumption,
throughput estimate, or batch-speedup result. For a fixed selection sequence
and fixed charge increments, changing only its elapsed execution speed does
not change the service-credit guarantee.

\subsection{Service metric and eligibility contract}
Let $\mathcal A$ be a \emph{fixed} cohort eligible at every selection boundary,
$U_i$ its cumulative analyzed service, and $D_0=\max_i U_i(0)-\min_i U_i(0)$.
An analyzed increment is nonnegative and at most $C_{\max}$. This metric may
be scheduler-billed service or useful protected service, but they must not be
interchanged. Actual synchronous callers can leave and rejoin eligibility;
the fixed-cohort assumption is therefore a refinement obligation, not a
synonym for high contention.

\subsection{Service-spread and accounting-error bounds}
\begin{theorem}[Approximate-min service bound]
\label{thm:spread}
If every selected client $i$ satisfies
$U_i\le\min_{j\in\mathcal A}U_j+\delta$ for fixed $\delta\ge0$, and no credit
is reset, then at every service boundary
\begin{equation}
 \max_i U_i-\min_i U_i\le B_U
   :=\max\{D_0,\delta+C_{\max}\}.\label{eq:spread}
\end{equation}
\end{theorem}
\begin{proof}
Assume the prior minimum is $m$ and spread at most $B_U$. Unchanged credits
remain in $[m,m+B_U]$. The selected credit becomes at most
$m+\delta+C_{\max}\le m+B_U$ and cannot decrease below $m$.
The new minimum cannot decrease, establishing the invariant by induction.
\end{proof}
Exact min gives $\delta=0$; a fixed eligibility window gives its window width.
For weights $w_i>0$, use $U_i=S_i/w_i$ and bound increments by
$\max_i c_{i,\max}/w_i$. This is a mathematical extension, not a claim that the
current FC-PQ implements weighted service or FC-EW.

\begin{corollary}[Bounded accounting error]
\label{cor:error}
Suppose the selector instead uses $V_i=U_i+e_i$ with
$\max_i e_i-\min_i e_i\le E_c$ at each selection, and chooses within $\delta$
of the minimum $V$. Equation~\eqref{eq:spread} holds with $\delta+E_c$ in
place of $\delta$.
\end{corollary}
\begin{proof}
For every $j$, $U_i+e_i\le U_j+e_j+\delta$, hence
$U_i\le U_j+\delta+E_c$. Apply Theorem~\ref{thm:spread}.
\end{proof}
A common credit offset is harmless. Small \emph{per-operation} timing errors
are not necessarily harmless: their cumulative pairwise discrepancy may
increase with unequal operation counts, so they do not establish a constant
$E_c$ for an unbounded execution.

\subsection{Billed service versus useful service}
\begin{proposition}[Billed fairness can bias useful service]
\label{prop:bias}
For fixed useful cost $c_i>0$ and fixed extra charge $\theta_i\ge0$, let each
operation add $q_i=c_i+\theta_i$ to billed credit. If all clients remain
backlogged, their cumulative billed credits diverge to infinity, and their
billed-credit differences remain bounded, their asymptotic useful-service
fractions are
\begin{equation}
 f_i=\frac{c_i/(c_i+\theta_i)}{\sum_j c_j/(c_j+\theta_j)}.\label{eq:bias}
\end{equation}
\end{proposition}
\begin{proof}
With $m_i$ completions, billed service is $m_iq_i$ and useful service is
$m_ic_i=(c_i/q_i)(m_iq_i)$. Bounded differences and divergence make the billed
service ratios tend to one. Normalize useful services by their sum.
\end{proof}
For $(c_1,c_2)=(1,9)$ and $\theta_1=\theta_2=1$, equal billed service gives
useful shares $(5/14,9/14)$, not $(1/2,1/2)$. This is an exact synthetic
example, not an estimated bias of the implementation. FC-PQ's outer timer
includes input/result handling and differs from the benchmark's inner
\texttt{hold\_time}; the distinction must be calibrated or explicitly chosen
as the fairness objective.

\paragraph{What JFI can establish.}
For positive mean cumulative service $\mu$ and variance $\sigma^2$, Jain's
index is $J=\mu^2/(\mu^2+\sigma^2)$. A range bound $B_U$ implies
$\sigma^2\le B_U^2/4$, and thus
$J\ge [1+B_U^2/(4\mu^2)]^{-1}$.
To see the variance bound, for values in $[m,M]$ expand
$(U-m)(M-U)\ge0$ and average to get
$\sigma^2\le(\mu-m)(M-\mu)\le(M-m)^2/4$.
The converse fails: $(U_1,U_2)=(k^2,k^2+k)$ has unbounded gap $k$ yet $J\to1$.
Windowed delivered service is a difference of cumulative values; if both
endpoints have bounded spread $B_U$, a per-pair window gap can be $2B_U$.
Do not attach this result to arbitrary newcomers or to a different charge
metric merely because a final JFI is near one.

\subsection{Waiting measured in competing services}
\begin{theorem}[Service volume before a visible request]
\label{thm:waiting}
Fix a target $i$ that remains eligible and unserved with credit $u=U_i$.
Keep the cohort fixed and the approximate-min condition of
Theorem~\ref{thm:spread}. If every service charge is at least $c_{\min}>0$,
the number of competing services before $i$ is selected is at most
\begin{equation}
 K_i=\sum_{j\ne i}\max\!\left\{0,
 \left\lfloor\frac{u+\delta-U_j}{c_{\min}}\right\rfloor+1\right\}.
 \label{eq:waiting}
\end{equation}
Their total charged service is at most
$\sum_{j\ne i}((u+\delta-U_j)_++C_{\max})$.
\end{theorem}
\begin{proof}
Before a competing service, approximate-min implies its current credit is
at most $u+\delta$, because $i$ remains eligible. After $k-1$ services of
$j$, its credit is at least $U_j+(k-1)c_{\min}$. Combining inequalities bounds
$k$ by the corresponding summand. Its final service can overshoot the
threshold by at most $C_{\max}$, giving the volume bound. Summing accounts
for every non-target selection; this also tolerates adversarial ties.
\end{proof}
If selection continues, the target must eventually be chosen. Zero-charge
operations without a tie-progress condition can defeat a selection-count
argument. This theorem concerns competing service, not wall-clock suspension
or request-response time. Translating it into elapsed time requires a
separate set of progress and timing assumptions.

\subsection{Applying the obligations to FC-PQ}
FC-PQ drains announcements once, then makes up to $H=64$ heap/tree pop
attempts. A completed entry consumes a pop but no service. It is temporarily
buffered or deactivated; active callers can republish without entering the
ring. Thus exact-min is with respect to the current selectable queue, not
automatically every logically pending requester. The code's attempted
starvation clamp is applied \emph{after} removing the minimum. Since the next
minimum is no smaller, $\min(\text{popped},\text{next})$ leaves credit unchanged.
The threshold provides no extra selection guarantee.

\paragraph{Admission counterexample.}
A zero-credit newcomer receives total charged service divided by total
completions: a mean \emph{per-operation} charge, not current cumulative virtual
time. After unit-cost history, an old pending client may have credit 1000
while the newcomer begins at 1. A repeatedly resubmitting newcomer remains
strictly preferred for 999 unit services until it ties the old credit.
Even the maximum 64 completions per pass requires at least 16 passes;
completed-node pops can require more. The post-pop threshold of eight cannot
promote the unpopped old client. This abstract schedule is consistent with
the selection rule; it is not a recorded Rust execution. Increasing the
history makes bypass arbitrarily large at fixed population and request cost.
It refutes a history-independent admission/response bound, not the correctly
scoped fixed-cohort theorem with its initial spread $D_0$.

\begin{proposition}[Announcement visibility in service units]
\label{prop:visibility}
Let $R\ge1$ and $H\ge0$ be integer budgets and $q\ge0$ a ticket count.
Suppose each FIFO drain consumes the first $\min(R,\text{queued tickets})$
reserved announcements, waiting for their publication if necessary, and at
most $H$ services occur between drains. Consider an announcement reserved
just after the current drain, with $q$ earlier unconsumed tickets ahead of it.
If drains continue and reserved tickets are eventually published/consumed
in order, the announcement is discovered after at most
\begin{equation}
 H\left\lceil\frac{q+1}{R}\right\rceil
 \label{eq:visibility}
\end{equation}
intervening service completions. This is not a wall-time guarantee.
\end{proposition}
\begin{proof}
Its ticket is in the $\lceil(q+1)/R\rceil$th future drain. There is at most
one $H$-service segment before each such drain, including the remainder of
the current pass. Later FIFO reservations cannot overtake the ticket.
\end{proof}
For FC-PQ, $R=H=64$ but these are different limits: ring admission versus PQ
pop attempts. The queue itself can contain more than 64 clients. A producer
reserves a slot before publishing its valid flag, and the combiner spins
until that flag appears. Arbitrarily delayed publication makes wall delay
unbounded even with few or zero intervening completions. Active-node reuse
and buffered reactivation are separate visibility paths, not covered by
Equation~\eqref{eq:visibility}. Together these observations explain why a
budget, visibility bound, usage bound, and response bound are distinct.

\subsection{Fairness of the implemented lock families}
\label{sec:fairness-instances}
Table~\ref{tab:fairness} compares selection and accounting contracts only.
The source evidence is shared with the implementation appendix, but no
communication cost or throughput ranking is used to infer these guarantees.

\begin{table*}[t]
\centering\small
\caption{Fairness comparison only. A policy description is not a proved
service bound; every conclusion is restricted to the stated metric,
eligible cohort, and progress assumptions.}
\label{tab:fairness}
\begin{tabular}{@{}P{0.14\textwidth}P{0.33\textwidth}P{0.47\textwidth}@{}}
\toprule
Family & Selection/admission contract & Service guarantee or missing obligation \\
\midrule
FC & Scan order over published active nodes & Scan order alone supplies no bounded cumulative usage gap for heterogeneous service costs. \\
FC-Ban / CC-Ban & Eligibility filtering at the combiner versus delay before enqueue & Neither penalty rule is proved here to implement approximate-min service. Define the eligible set and charge metric separately for each path. \\
FC-PQ heap/tree & Minimum recorded credit in the currently selectable PQ & Fixed-cohort theorem is conditional on continuous selectability and bounded charges. Newcomer initialization, reactivation, and the inert clamp prevent an unconditional dynamic-arrival bound. \\
FC-SL & Ordered SkipSet with distinct publication and charging boundaries & Do not transfer FC-PQ's admission proof. Shared non-atomic counters are a correctness blocker, not a bounded accounting error. \\
CC / DSM & Linked FIFO combining & FIFO service order is not equal cumulative service. Progress through unpublished links is an additional obligation. \\
MCS / CLH / Ticket & Ordered queue or ticket acquisition & With continuing progress, acquisition order still does not equalize heterogeneous held service. No min-credit theorem follows from FIFO alone. \\
Spin / Mutex & Barging or runtime-specific acquisition & No usage-fair selection theorem is supplied by these interfaces; runtime progress must be specified independently. \\
Shfl / CFL & Queue reshaping; CFL accounts per core/NUMA node & Locality order and accounting domains differ from per-requester service. Do not transfer an FC-PQ $C_{\max}$ bound by analogy. \\
U-SCL & Slice admission with held-time and weight-based billing & Its actual admission and integer-weight charge require a separate metric/progress argument, not inheritance of the PQ theorem. \\
\bottomrule
\end{tabular}
\end{table*}

\paragraph{Acquisition order is not usage fairness.}
For fixed useful costs $c_i,c_j$, a cyclic schedule serving each client once
per round yields a cumulative service gap $r|c_i-c_j|$ after $r$ rounds from
equal initial service. Thus even repeated equal acquisition counts need
not bound usage disparity. This argument uses no machine-speed assumption.

\subsection{Fairness validation and limits}
\label{sec:fairness-validation}
The artifact's \src{spread}, \src{accounting_error},
\src{waiting_bound}, \src{charge_bias}, and
\src{announcement_visibility} checks exercise the stated invariants and
service-count arguments. Counterexamples challenge dynamic-admission,
clamp, zero-charge progress, and final-JFI claims. These finite checks
supplement the proofs; none certifies an executed Rust trace or memory safety.

Implementation refinement must track invocation, publication and selection
visibility, cohort entry/exit, initial credit, each billed increment, and
useful protected service separately. Test cumulative and windowed service
gaps, rather than treating high final JFI or high throughput as sufficient.
The fixed-additive-charge model predicts unequal useful-service fractions
despite bounded billed-credit spread; its magnitude in the implementation
depends on the actual charge boundary, not on a synthetic numerical example.
No hardware cost calibration establishes continuous eligibility or correct
accounting. Dynamic membership needs its own service contract before any
fixed-cohort theorem can be applied.

\clearpage
\balance

\section{Discussion: relating the two analyses}
\label{sec:discussion}
\paragraph{Independent outputs.}
The performance analysis yields source-work demand, throughput conditions,
and return-tail certificates for a supplied workload and execution policy.
The fairness analysis yields service-spread and competing-selection bounds,
plus the admission and accounting conditions under which they hold.
Low amortized cost proves no service guarantee; bounded credit spread proves
no speedup, cache residency, or wall-clock deadline.

\paragraph{An explicit interface, not a circular argument.}
A policy can change the executed class mix, actual batch occupancy, and
executor sequence. The performance analysis consumes those quantities as
measured or separately justified inputs. For example, useful-service
fractions from Proposition~\ref{prop:bias} can be substituted into
Equation~\eqref{eq:mix} only when both sets of assumptions hold.
Its billing overhead is a credit increment, not automatically another
elapsed-time term to add to the performance equation.

Similarly, a fairness bound on competing selections is not a latency bound.
Conversion additionally requires bounded admission/discovery, acquisition,
per-selection elapsed cost, preemption, completion observation, and any
post-completion combining work. The performance return-tail envelope covers
only that last component under its stated assumptions. No unconditional
$O(C_{\max})$ request-response bound follows by combining the two analyses.

\paragraph{Evaluating a design tradeoff.}
Compare implementation costs at controlled work and placement, then report
the service guarantee separately. Intra-epoch executor invariance explains
why reordering requests need not introduce protected-state handoffs, but
does not remove scheduling cost or establish that the ordering is fair.
FC-PQ versus FC and CFL versus MCS must therefore be reported as pairs of
performance and fairness conclusions, not as a single ``fairness is free''
claim.

\paragraph{Evidence boundary.}
\label{sec:evidence}
The same executable artifact runs two independent groups of mathematical
checks: crossover/return feasibility for performance, and service spread,
accounting error, waiting, charging bias, and visibility for fairness.
It also emits explicit counterexamples. Written arguments establish the
conditional results; finite enumeration supplements them. Neither group is
an executed Rust schedule, a weak-memory model check, or a machine benchmark.
The shared appendix records implementation facts and reproduction without
conflating the two kinds of guarantee.

\begin{thebibliography}{9}
\bibitem{logp} D. Culler et al. LogP: Towards a Realistic Model of Parallel Computation. PPoPP, 1993. \url{https://doi.org/10.1145/155332.155333}.
\bibitem{loggp} A. Alexandrov, M. F. Ionescu, K. E. Schauser, and C. Scheiman. LogGP. SPAA, 1995. \url{https://doi.org/10.1145/215399.215427}.
\bibitem{coherence} S. Ramos and T. Hoefler. Modeling Communication in Cache-Coherent SMP Systems---A Case-Study with Xeon Phi. HPDC, 2013. \url{https://spcl.inf.ethz.ch/Publications/index.php?pub=162}.
\bibitem{fc} D. Hendler, I. Incze, N. Shavit, and M. Tzafrir. Flat Combining and the Synchronization-Parallelism Tradeoff. SPAA, 2010, pp. 355--364.
\bibitem{cc} P. Fatourou and N. D. Kallimanis. Revisiting the Combining Synchronization Technique. PPoPP, 2012.
\bibitem{mcs} J. M. Mellor-Crummey and M. L. Scott. Algorithms for Scalable Synchronization on Shared-Memory Multiprocessors. ACM TOCS, 1991. \url{https://doi.org/10.1145/103727.103729}.
\end{thebibliography}
\clearpage
\nobalance
\onecolumn
\appendix
\section{Implementation inventory and evidence}
\label{app:inventory}
This appendix records source deductions at implementation commit
\texttt{56a02ab}; analysis files are additional worktree changes. It is a static
protocol audit, not a build, concurrency test, or assertion that every path
is correct. The inventory is determined by
\src{src/command_parser/lock_target.rs}:110--154,183--215 and
\src{crates/libdlock/src/dlock2.rs}:55--81, rather than older README lists.
There are 21 DLock2 target variants, including the two PQ backends and both
Rust/C shuffle and CFL implementations. Similar names do not establish
identical execution, timing, or fairness contracts.
This is a shared factual inventory, not a combined analytical model.
Performance conclusions are in Section~\ref{sec:performance}; fairness
conclusions are in Section~\ref{sec:fairness}, each with its own comparison.

For compact source references below, \src{d2/} means
\src{crates/libdlock/src/dlock2/}, \src{d1/} means
\src{crates/libdlock/src/dlock/}, and \src{lib/} means
\src{crates/libdlock/src/}. Paths beginning \src{c/} or \src{src/} are
relative to the repository root. Line ranges refer to this baseline.

\subsection{All DLock2 targets}
\small
\begin{longtable}{@{}P{0.14\textwidth}P{0.51\textwidth}P{0.27\textwidth}@{}}
\toprule
Target enum & Model instance, budget, and guarantee boundary & Source evidence \\
\midrule
\endfirsthead
\toprule Target enum & Model instance, budget, and guarantee boundary & Source evidence \\
\midrule
\endhead
\bottomrule
\endfoot
FC & Inactive nodes publish through a head CAS; each pass scans the retained list to its end. No count/time cap. Cleanups run at pass multiples of 500 and remove sufficiently old nodes. Demand is list visits per completion, with election/activation retries separate; scan order does not establish usage fairness. & \src{d2/fc/lock.rs}:17,53--138,152--190. \\
\addlinespace
FCBan & Full FC-style scan plus combiner-side deadline tests. A skip does not complete a request. Penalties use elapsed work intervals and a counter of retained registrations, not instantaneous pending clients; eligibility and exposed server idle must be separated. & \src{d2/fc_ban/lock.rs}:59--62,103--138,148--167,182--205. \\
\addlinespace
CC & One tail swap per call, input/link publication in the predecessor slot, completion/wait flags, and baton handoff. At most 64 executed delegates per batch, possibly fewer on missing links. FIFO request order is not equal usage. & \src{d2/cc/lock.rs}:23,52--86,97--121. \\
\addlinespace
CCBan & TSC/backoff admission before enqueue, then queue combining. Cap 16 rather than CC's 64. Charge uses work intervals times a registration counter, which increments on first thread-local initialization. Thus the CC--CCBan delta is not a pure fairness-policy ablation. & \src{d2/cc_ban/lock.rs}:50--96,145--184. \\
\addlinespace
DSM & Per-client two-node toggle, tail enqueue, linked FIFO service. The post-execution test \texttt{counter > H} with default 64 permits 65 callbacks, subject to earlier successor/lookahead exits. It is not global double-buffer rotation. Unpublished successors can delay handoff. & \src{d2/dsm/lock.rs}:27,69--101,112--176. \\
\addlinespace
FcSL & Callers insert into a shared ordered SkipSet; a pass tries at most 64 pops and checks completion before executing. Successful service deactivates the node; a new call republishes. Charges include time since the previous timestamp, not only the delegate. Newcomer mean-charge initialization differs from a cumulative baseline. & \src{d2/fc_sl/lock.rs}:72--99,102--150,161--195. \\
\addlinespace
FcPqBTree & Shared sequential BTreeSet under combiner ownership, with a 64-slot announcement ring and at most 64 PQ pops. Ordinary queue work scales with arrivals/pops and queue size; allocation and migrating metadata matter. Exact-min proof is conditional on fixed continuous selectability and the billed metric. & \src{d2/fc_pq/lock.rs}:109--255,275--313; \src{lib/sequential_priority_queue.rs}:39--61. \\
\addlinespace
FcPqBHeap & Same admission, service policy and buffers as the tree version; BinaryHeap reverses ordering to pop minimum usage. Heap growth can relocate storage. Backend locality costs cannot be inferred from the identical min-credit policy. & \src{d2/fc_pq/lock.rs}:145--255; \src{lib/sequential_priority_queue.rs}:14--37. \\
\addlinespace
SpinLock & Contended acquisition operates on one AtomicBool with backoff; release stores false. The wrapper executes one delegate on the caller. Count attempts separately from ownership transfers; no FIFO or usage guarantee is supplied. & \src{lib/spin_lock.rs}:30--50; \src{d2/spinlock.rs}:41--46. \\
\addlinespace
MCS & Tail swap, predecessor link publication, own-node polling, and successor flag handoff or empty-tail CAS. One caller-executed operation per epoch. Constant successful queue operations do not bound a delayed link publication or protected-data migration latency. & \src{d2/mcs.rs}:78--107,124--152; \src{d2/spinlock.rs}:41--46. \\
\addlinespace
CLH & Tail swap, predecessor-flag polling, own-node release, and predecessor-node recycling. FIFO acquisition, not equal service. Polling is distributed over nodes but enqueue is a shared hotspot. & \src{d2/clh.rs}:98--124,171--185. \\
\addlinespace
Ticket & Fetch-add ticket admission and shared \texttt{now\_serving} polling; release advances the turn. Logical sharer invalidation demand is distinct from message count and latency. FIFO turn order does not bound heterogeneous service spread. & \src{d2/ticket.rs}:37--53. \\
\addlinespace
Mutex & One std::sync::Mutex guard and one delegate. The adapter supplies no fairness policy and does not pin a standard-library/kernel waiting implementation. Its opaque costs require the actual runtime and OS, not a guessed futex constant. & \src{d2/mutex.rs}:21--24,36--44. \\
\addlinespace
PthreadMutex & Explicit PTHREAD\_MUTEX\_NORMAL, lock/delegate/unlock. Libc/kernel contention, wakeup and scheduling remain external parameters; the type name does not supply an admission fairness bound. & \src{d2/pthread_mutex.rs}:60--65,89--100. \\
\addlinespace
ShflLock & Caller execution with fast-path ownership, queue promotion and same-NUMA shuffling. Leader traversal can stop at queue boundaries or a PRNG condition; no fixed request batch. \texttt{wcount} is not a delegate count budget. & \src{d2/shfl_lock/lock.rs}:23--30,200--289,298--417. \\
\addlinespace
ShflLockC & RawCAqs FFI reaches the Cargo-built AQS source. Same placement-policy family, but count its own traversal, promotion and lock-byte events; one delegate per acquisition, not combining. & \src{d2/c_aqs.rs}:75--88; \src{c/shfllock/aqs.c}:122--184,220--313. \\
\addlinespace
CFL & Queue shuffling plus global per-core/per-NUMA runtime accounting; a node-runtime scan and variable queue traversal affect policy demand. The fairness entity is not automatically an FC-PQ requester. Executor topology changes both policy work and $M_H$. & \src{d2/cfl.rs}:354--390,397--534,554--684,721--745. \\
\addlinespace
CFLC & RawCCfl FFI reaches the C implementation's runtime-based queue reshaping. One delegate per hold. Global core/node counters and a 100,000-tick node-spread threshold are source policy inputs, not a per-thread $C_{\max}$ theorem. Rust CFL also uses core/node accounting; neither should be described as an ideal per-request virtual-time scheduler. & \src{d2/c_cfl.rs}:75--90; \src{c/cfl/cfl.c}:187--304,312--445. \\
\addlinespace
FcC & Full active-list scan; completion is the delegate-pointer release to null after the response write. Cleanup uses 50-pass aging rather than Rust FC's 500. Waiters yield/retry after 100 pause iterations. The active scan has no enforced time cap merely because the build defines a timing macro. & \src{lib/c_binding/flatcombining.rs}:75--131; \src{c/FlatCombining/original/flatcombining.c}:13--56,90--129. \\
\addlinespace
CcC & FIFO tail queue, predecessor wait/completed flags, at most 64 callbacks while a successor exists, and a trailing-node combiner baton. This cap is established in the C source, not transferred from the Rust name. & \src{lib/c_binding/ccsynch.rs}:78--134; \src{c/CCsynch/ccsynch.c}:39--88. \\
\addlinespace
USCL & One caller-held critical section with slice admission, spin/yield, futex wait/wake and banned-time sleep. The nominal 2\,ms slice uses compile-time cycle conversion. Charge is held TSC ticks times an integer weight quotient, not delegate-only time. Default lazy C weighting is not an explicit DLock2 requester-weight policy. & \src{d2/uscl.rs}:21--41; \src{c/u-scl/fairlock.h}:91--185,188--305. \\
\end{longtable}
\normalsize

\paragraph{The actual C build.}
\src{crates/libdlock/build.rs}:18--33 compiles
\src{c/CCsynch/ccsynch.c},
\src{c/FlatCombining/original/flatcombining.c},
\src{c/u-scl/fairlock.c}, \src{c/shfllock/aqs.c}, and \src{c/cfl/cfl.c}.
\src{c/u-scl/fairlock.c}:1 includes the header containing its implementation.
The build defines \texttt{CYCLE\_PER\_US=2400}; this is not a measurement of
invariant-TSC frequency. A defined \texttt{FC\_THREAD\_MAX\_CYCLE} does not
establish a time bound in the original FC scan, which traverses to null.
Binding headers are in \src{crates/libdlock/binding/wrapper.h}:1--5.
The raw-lock DLock2 wrapper executes exactly one delegate between acquire and
release (\src{d2/spinlock.rs}:41--46), so these wrappers do not gain delegation
batching from sharing the DLock2 interface.

\subsection{FC-PQ proof and measurement obligations}
\begin{enumerate}
\item \textbf{Admission and visibility.}
\src{d2/fc_pq/lock/buffer.rs}:52--84 reserves by tail fetch-add, waits for
capacity, then publishes a valid flag. The iterator snapshots at most 64
tickets and spins on each valid flag before consuming it (:96--150).
This supports Proposition~\ref{prop:visibility}'s FIFO service-count bound,
not a wall-clock deadline. Counter wraparound and memory correctness are
excluded from that abstraction.
\item \textbf{Selectable versus merely pending.}
\src{d2/fc_pq/lock.rs}:175--255 buffers completed nodes, removes inactive ones,
and may reinsert reactivated nodes. A request becoming ready while absent
from the PQ need not be selected at the same boundary as a globally visible
minimum. Newcomer credit is the mean operation charge (:157--162).
No uniform admission-independent initial-spread bound follows.
\item \textbf{Attempted starvation prevention.}
The supported heap/tree backends pop a minimum
(\src{lib/sequential_priority_queue.rs}:14--61). Before the next queue mutation,
the post-pop clamp at \src{d2/fc_pq/lock.rs}:193--197 cannot decrease that
minimum. Its threshold cannot promote an entry that has not been popped.
\item \textbf{Billing and return.}
The outer timestamps bracket input access, delegate call and result write
(:209--217). The benchmark's inner useful-work timer has different boundaries
(\src{src/benchmark/dlock2/proportional_counter.rs}:81--110).
A combiner unlocks before checking its own completion (:286--295), so an
operation's finish and the caller's return are distinct.
\item \textbf{Misleading timer boundary.}
The same \texttt{begin} variable is set at pass entry and reset at every
delegate before \texttt{combiner\_time\_stat += end-begin}
(:129--131,209,258--263). This does not accumulate the entire pass when it
serves requests. FC-SL similarly resets its timestamp after each service
(\src{d2/fc_sl/lock.rs}:102--150). These statistics cannot be substituted for
full-pass occupancy without correcting the instrumentation.
\end{enumerate}
These are source-derived obligations and counterexamples, not silently fixed
production algorithms. In particular, proving an ideal policy's invariant
does not certify an unsafe concurrent implementation.

\subsection{FC-SL: concrete synchronization prerequisite}
\label{app:fcsl-race}
The DLock2 FC-SL newcomer path has a source-level data race. Its
\texttt{total\_usage} and \texttt{total\_served} fields are non-atomic
\texttt{SyncUnsafeCell<u64>} (\src{d2/fc_sl/lock.rs}:48,50).
Every \texttt{push\_node} reads the served counter, and zero-usage initialization
also reads total usage (:78--82). This path runs from \texttt{lock()} before
attempting the combiner mutex (:169--174), while another thread holding that
mutex can update both counters in \texttt{combine()} (:134--136).
There is no common synchronization protecting these conflicting accesses.
The per-request active/completion flags and a later SkipSet insertion do not
protect the global counters. Interior mutability and aligned x86 loads do
not legalize a Rust data race.

This is a correctness prerequisite for the concrete FC-SL instantiation,
not a claim about the magnitude of accounting error $E_c$. Moving initialization
under combiner ownership or adding synchronized counters with an explicit
snapshot-consistency contract are possible repairs, but no repair is made in
this analysis worktree. No executed race detector result or complete
concurrency audit is claimed; the conflicting source paths are sufficient
to identify this missing synchronization. Other variants have not thereby
been certified race-free.

\subsection{Legacy DLock1 coverage}
The 11 target variants are listed in
\src{src/command_parser/lock_target.rs}:41--66 and constructed at :83--105.
DLock1 passes callbacks over a guard; it is not the same API/charge boundary as
DLock2's function-delegate input/result path. The table accounts for every
legacy target, rather than transferring a DLock2 bound by matching its name.

\small
\begin{longtable}{@{}P{0.14\textwidth}P{0.51\textwidth}P{0.27\textwidth}@{}}
\toprule
Target & Instance and distinguishing source behavior & Evidence \\
\midrule
\endfirsthead
\toprule Target & Instance and distinguishing source behavior & Evidence \\
\midrule
\endhead
\bottomrule
\endfoot
FC & Callback-based full linked-list scan, with Parker wake/completion events; no fixed count/time cutoff in the scan. & \src{d1/fc/fclock.rs}:71--100. \\
\addlinespace
FCBan & Combiner-side banned-until filter over the retained list. Full scan and eligibility tests remain; do not transplant DLock2's exact penalty formula without checking its charge/update code. & \src{d1/fc_fair_ban.rs}:77--123,165--203. \\
\addlinespace
FCBanSlice & Despite its name and declared slice constant, the early-return slice test is commented out. The active combiner traverses to list end. Accumulated work affects an adaptive caller wait timeout, not an enforced pass deadline. & \src{d1/fc_fair_ban_slice.rs}:82--154,210--224. \\
\addlinespace
CC & Current combiner request is included in the \texttt{H=16} loop. Up to 16 callbacks before baton handoff; different from DLock2 CC's 64. & \src{d1/ccsynch.rs}:104--132. \\
\addlinespace
CCBan & Separately executes the current request, then up to 16 followers: up to 17 callbacks per episode. Adds pre-enqueue banned-time delay. Thus its total budget differs from DLock1 CC despite the same loop constant 16. & \src{d1/ccsynch_fair_ban.rs}:61--75,99--118,153--197. \\
\addlinespace
FCSLNaive & Usage-keyed concurrent SkipMap, not a different non-concurrent container inferred from ``Naive.'' Nominal accumulated work threshold 240,000 TSC units, checked after an item; callers retry after timed waiting without an explicit combiner baton. & \src{d1/fc_sl_naive.rs}:22--40,57--110,119--153. \\
\addlinespace
FCSL & Same SkipMap/threshold but adds completion tracking and wakes a remaining queued node as next combiner. Charging includes dequeue and surrounding work. Neither its post-item threshold nor calibrated cycle convention establishes preemptive execution. & \src{d1/fc_sl.rs}:22--40,57--126,134--166. \\
\addlinespace
RCL & Dedicated server thread scans and executes posted callbacks. Generic conversion intentionally returns None, but a separate benchmark path constructs it. Server-resource accounting is essential; this target is not absent from DLock1. & \src{src/benchmark/dlock.rs}:107--145; \src{d1/rcl/rclserver.rs}:85--98; \src{d1/rcl/rclthread.rs}:42--97. \\
\addlinespace
Mutex & Non-delegation OtherLocks baseline; uses the mutex callback adapter, not a parameterized Parker. Runtime/OS waiting remains a separate model component. & \src{src/command_parser/lock_target.rs}:98--102; \src{lib/dlock.rs}:55--119. \\
\addlinespace
SpinLock & Non-delegation OtherLocks baseline, not converted into a futex-Parker implementation by selecting the block waiter option. & \src{src/command_parser/lock_target.rs}:98--102; \src{src/benchmark/dlock.rs}:38--57. \\
\addlinespace
USCL & OtherLocks wrapper over the scheduler-cooperative acquisition protocol; no DLock1 Parker parameter is applied. Slice/OS costs must be modeled from its underlying implementation. & \src{src/command_parser/lock_target.rs}:98--102; \src{lib/u_scl.rs}:41--51. \\
\end{longtable}
\normalsize

\paragraph{Parking is a model dimension.}
For delegation targets, the default \texttt{All} waiter choice expands into
SpinParker and BlockParker variants
(\src{src/benchmark/dlock.rs}:38--57). SpinParker uses polling/backoff and may
yield (\src{lib/parker/spin_parker.rs}:26--63); BlockParker adds futex wait and
wake events (\src{lib/parker/block_parker.rs}:18--100). Equal service-credit
bounds do not remove scheduler wakeup delays. Baselines in OtherLocks do not
inherit this Parker selection.

\paragraph{RCL changes the resource budget, not only migration.}
\src{src/benchmark/dlock.rs}:125--145 starts the server using
\texttt{num\_cpu-1} and supplies that reduced CPU count to client work while
retaining the thread count. The server thread requests affinity at
\src{d1/rcl/rclthread.rs}:42--97. Assuming successful pinning, a single-server
run can eliminate rotating-executor migration after warmup, but still pays
request publication/completion and scanning. Its server is a resource in
Proposition~\ref{prop:resource}, and worker capacity is reduced; an apples-to-
apples comparison must hold total hardware resources fixed. This is source
reachability and requested placement, not a measurement that affinity always
succeeds. The exported DLock2 \src{d2/rcl.rs} modules have no DLock2 enum/CLI
variant (\src{lib/dlock2.rs}:25--26,55--81); direct use is not a dispatched
DLock2 benchmark and is not asserted working here.

\section{Reproduction and claim/evidence boundary}
\label{app:artifact}
The standalone manuscript is \src{analysis/logp/paper.tex}.
\src{performance.tex} and \src{fairness.tex} contain separate analyses;
\src{section.tex} supplies their common scope and final discussion.
The dependency-free artifact is \src{analysis/logp/check_model.py},
requiring Python 3.9 or newer. From the
worktree root:
\begin{quote}\ttfamily\small
python3 analysis/logp/check\_model.py\par
\quad --output .worktree/logp/verification.json
\end{quote}
The line break above is typographical: run it as one command. The adjacent
README supplies copyable commands for the checker and \texttt{latexmk} build.
Outputs are in the ignored \src{.worktree/logp/} directory. All numerical
illustrations in the section are exact normalized model values, not machine
measurements.

\begin{center}\small
\begin{tabular}{lrl}
\toprule
Check & Cases & Enumerated evidence \\
\midrule
\multicolumn{3}{l}{\emph{Performance}}\\
Crossover &5,120& Rational inequality cases \\
Batch frontier &5,832& Brute-force versus algebraic integer intervals \\
\addlinespace
\multicolumn{3}{l}{\emph{Fairness}}\\
Spread &44,400& One-step normalized integer invariant transitions \\
Accounting error &164,421& True/virtual-credit error transitions \\
Waiting bound &22,587& Competing transitions and terminal target choices \\
Charge bias &256& Complete equal-billed-service cycles \\
Announcement visibility &1,485& FIFO drain/service schedules \\
\addlinespace
\multicolumn{3}{l}{\emph{Cross-checks}}\\
Counterexamples &7& Source-inspired and metric witness families \\
\bottomrule
\end{tabular}
\end{center}
All eight checks passed in both standard and optimized (\texttt{python3 -O})
executions; the generated JSON reports were byte-identical. Their case counts
have different meanings and are not numbers of independent concurrent Rust
tests. Parameter domains are explicit in the script. It enumerates all
permitted choices within those domains rather than sampling random traces;
written proofs, not the finite domains, establish the conditional unbounded
results. The script uses explicit failures rather than optimization-removable
\texttt{assert} statements. Its JSON includes assumptions and example witnesses.

\begin{center}\small
\begin{tabular}{@{}P{0.22\textwidth}P{0.33\textwidth}P{0.36\textwidth}@{}}
\toprule
Claim & Evidence supplied & What is not inferred \\
\midrule
Abstract resource/migration bounds & Stated machine/residency assumptions and proofs & Guaranteed physical costs from empirical averages \\
Batch crossover and return frontier & Algebra, integer certificates, exact-rational checks & A hardware speedup or unconditional FC-PQ return deadline \\
Fixed-cohort service bounds & Inductive proofs, error reduction, finite transitions & Fairness under arbitrary admissions, suspension or credit resets \\
Charging bias & Exact fixed-cost derivation and complete-cycle checks & The magnitude of bias in actual benchmark data \\
FC-PQ visibility/newcomer/clamp limits & Source paths and abstract counterexamples & Executed Rust adversarial traces or full concurrency correctness \\
Coverage of targets & Enum/dispatch, budget and charge source audit & Successful compilation, safety, or performance of every variant \\
\bottomrule
\end{tabular}
\end{center}

The prescribed development environment previously failed evaluation because
\texttt{dotenv.resolved} had no defined value. The standard-library checks
and manuscript tooling use existing host programs instead; environment files
and lock implementations are unchanged. No performance suite is reported.
A full systems paper still needs independent calibration, corrected
measurement boundaries, controlled experimental resource accounting and
held-out prediction. These empirical requirements do not turn the proven
conditional statements into mere proposals, and the conditional statements
do not substitute for the empirical requirements.

\end{document}
